% ============================================================
% LATEX_GUIDE.tex
% A practical LaTeX reference for CS thesis students
%
% HOW TO USE THIS FILE:
%   This file is NOT included in your compiled thesis.
%   It is a reference you can read alongside your chapter files.
%   Every example below is valid LaTeX — copy and paste directly
%   into your chapter files and edit as needed.
%
%   To include this file in your compiled output temporarily
%   (e.g., to see how examples render), add this line to
%   Thesis.tex just before \end{document}:
%     % ============================================================
% LATEX_GUIDE.tex
% A practical LaTeX reference for CS thesis students
%
% HOW TO USE THIS FILE:
%   This file is NOT included in your compiled thesis.
%   It is a reference you can read alongside your chapter files.
%   Every example below is valid LaTeX — copy and paste directly
%   into your chapter files and edit as needed.
%
%   To include this file in your compiled output temporarily
%   (e.g., to see how examples render), add this line to
%   Thesis.tex just before \end{document}:
%     % ============================================================
% LATEX_GUIDE.tex
% A practical LaTeX reference for CS thesis students
%
% HOW TO USE THIS FILE:
%   This file is NOT included in your compiled thesis.
%   It is a reference you can read alongside your chapter files.
%   Every example below is valid LaTeX — copy and paste directly
%   into your chapter files and edit as needed.
%
%   To include this file in your compiled output temporarily
%   (e.g., to see how examples render), add this line to
%   Thesis.tex just before \end{document}:
%     % ============================================================
% LATEX_GUIDE.tex
% A practical LaTeX reference for CS thesis students
%
% HOW TO USE THIS FILE:
%   This file is NOT included in your compiled thesis.
%   It is a reference you can read alongside your chapter files.
%   Every example below is valid LaTeX — copy and paste directly
%   into your chapter files and edit as needed.
%
%   To include this file in your compiled output temporarily
%   (e.g., to see how examples render), add this line to
%   Thesis.tex just before \end{document}:
%     \include{LATEX_GUIDE}
%   Remove it before your final submission.
% ============================================================

\chapter*{LaTeX Quick Reference}
\addcontentsline{toc}{chapter}{LaTeX Quick Reference}

\noindent
This chapter is a reference for common LaTeX features used in
CS theses. It is intended for students who are new to LaTeX or
who need a quick reminder of the syntax for specific tasks.

% ============================================================
\section*{1. Basic Text Formatting}
% ============================================================

\noindent
LaTeX handles most formatting automatically. A few things you
will use frequently:

\begin{itemize}
    \item \textbf{Bold text}: \verb|\textbf{your text}|
    \item \textit{Italic text}: \verb|\textit{your text}|
    \item \texttt{Code inline}: \verb|\texttt{your text}| or \verb|\verb|CODE||
    \item Line break within a paragraph: \verb|\\|
    \item New paragraph: leave a blank line in the source
    \item Em dash (—): \verb|---| \quad En dash (–): \verb|--|
    \item Quotation marks: \verb|``text''| renders as ``text''
\end{itemize}


% ============================================================
\section*{2. Lists}
% ============================================================

\noindent
\textbf{Bulleted list:}

\begin{verbatim}
\begin{itemize}
    \item First item
    \item Second item
    \item Third item
\end{itemize}
\end{verbatim}

Renders as:
\begin{itemize}
    \item First item
    \item Second item
    \item Third item
\end{itemize}

\noindent
\textbf{Numbered list:}

\begin{verbatim}
\begin{enumerate}
    \item First item
    \item Second item
\end{enumerate}
\end{verbatim}

\noindent
\textbf{Custom labels} (for research questions, contributions, etc.):

\begin{verbatim}
\begin{enumerate}
    \item[\textbf{RQ1}] First research question.
    \item[\textbf{RQ2}] Second research question.
\end{enumerate}
\end{verbatim}


% ============================================================
\section*{3. Figures}
% ============================================================

\noindent
\textbf{Single figure:}

\begin{verbatim}
\begin{figure}[H]
    \centering
    \includegraphics[width=0.7\textwidth]{Images/your_figure.png}
    \caption{A descriptive caption explaining the figure.}
    \label{fig:your_label}
\end{figure}
\end{verbatim}

\noindent
Key options for \verb|\includegraphics|:
\begin{itemize}
    \item \verb|width=0.7\textwidth| — 70\% of text width (adjust as needed)
    \item \verb|width=\textwidth| — full text width
    \item \verb|height=5cm| — fixed height (width scales automatically)
    \item \verb|scale=0.5| — scale to 50\% of original size
\end{itemize}

\noindent
\textbf{Reference a figure in your text:}

\begin{verbatim}
Figure~\ref{fig:your_label} shows...
As shown in Figure~\ref{fig:your_label},...
\end{verbatim}

\noindent
The \verb|~| prevents a line break between ``Figure'' and the number.

\noindent
\textbf{Side-by-side figures:}

\begin{verbatim}
\begin{figure}[H]
    \centering
    \begin{subfigure}[t]{0.48\textwidth}
        \centering
        \includegraphics[width=\textwidth]{Images/figure_a.png}
        \caption{Caption for the left figure.}
        \label{fig:left}
    \end{subfigure}
    \hfill
    \begin{subfigure}[t]{0.48\textwidth}
        \centering
        \includegraphics[width=\textwidth]{Images/figure_b.png}
        \caption{Caption for the right figure.}
        \label{fig:right}
    \end{subfigure}
    \caption{Overall caption for both figures.}
    \label{fig:both}
\end{figure}
\end{verbatim}

\noindent
\textbf{Important:} Put all image files in the \texttt{Images/} folder.
Supported formats: PNG, JPG, PDF. PNG is recommended for screenshots
and diagrams. PDF is recommended for vector graphics.


% ============================================================
\section*{4. Tables}
% ============================================================

\noindent
\textbf{Basic table:}

\begin{verbatim}
\begin{table}[H]
    \centering
    \caption{Caption goes ABOVE the table.}
    \label{tab:your_label}
    \begin{tabular}{lcc}
        \toprule
        \textbf{Column 1} & \textbf{Column 2} & \textbf{Column 3} \\
        \midrule
        Row 1 data        & 0.92              & 0.88              \\
        Row 2 data        & 0.87              & 0.91              \\
        Row 3 data        & 0.95              & 0.93              \\
        \bottomrule
    \end{tabular}
\end{table}
\end{verbatim}

\noindent
Column alignment characters:
\begin{itemize}
    \item \texttt{l} — left aligned
    \item \texttt{c} — centered
    \item \texttt{r} — right aligned
    \item \texttt{p\{3cm\}} — fixed-width column with text wrapping
\end{itemize}

\noindent
\textbf{Reference a table:}

\begin{verbatim}
Table~\ref{tab:your_label} shows...
\end{verbatim}

\noindent
\textbf{Auto-width table} (fills text width, good for many columns):

\begin{verbatim}
\begin{table}[H]
    \centering
    \caption{Caption here.}
    \label{tab:wide}
    \begin{tabularx}{\textwidth}{lXX}
        \toprule
        \textbf{Method} & \textbf{Description} & \textbf{Result} \\
        \midrule
        Method A & Describe it here & 92.3\% \\
        Method B & Describe it here & 88.7\% \\
        \bottomrule
    \end{tabularx}
\end{table}
\end{verbatim}

\noindent
\texttt{X} columns expand to fill available space equally.


% ============================================================
\section*{5. Equations}
% ============================================================

\noindent
\textbf{Numbered equation:}

\begin{verbatim}
\begin{equation}
    \text{Accuracy} = \frac{TP + TN}{TP + TN + FP + FN}
    \label{eq:accuracy}
\end{equation}
\end{verbatim}

\noindent
Renders as:
\begin{equation}
    \text{Accuracy} = \frac{TP + TN}{TP + TN + FP + FN}
    \label{eq:accuracy}
\end{equation}

\noindent
\textbf{Reference it in text:}

\begin{verbatim}
Equation~\ref{eq:accuracy} defines...
\end{verbatim}

\noindent
\textbf{Inline math} (within a sentence):

\begin{verbatim}
The learning rate $\alpha = 0.001$ was fixed throughout training.
\end{verbatim}

\noindent
\textbf{Unnumbered equation:}

\begin{verbatim}
\[
    y = mx + b
\]
\end{verbatim}

\noindent
\textbf{Multi-line aligned equations:}

\begin{verbatim}
\begin{align}
    f(x) &= x^2 + 2x + 1 \\
         &= (x + 1)^2
    \label{eq:multi}
\end{align}
\end{verbatim}

\noindent
Common math symbols:
\begin{itemize}
    \item Fraction: \verb|\frac{numerator}{denominator}|
    \item Subscript: \verb|x_i| → $x_i$
    \item Superscript: \verb|x^2| → $x^2$
    \item Greek: \verb|\alpha \beta \gamma \sigma \mu| → $\alpha\ \beta\ \gamma\ \sigma\ \mu$
    \item Sum: \verb|\sum_{i=1}^{n}| → $\sum_{i=1}^{n}$
    \item Square root: \verb|\sqrt{x}| → $\sqrt{x}$
    \item Infinity: \verb|\infty| → $\infty$
    \item Approximately: \verb|\approx| → $\approx$
    \item Less/greater or equal: \verb|\leq \geq| → $\leq\ \geq$
\end{itemize}


% ============================================================
\section*{6. Code Listings}
% ============================================================

\noindent
\textbf{Inline code} (short snippets in a sentence):

\begin{verbatim}
The function \texttt{train\_model()} accepts a learning rate parameter.
\end{verbatim}

\noindent
\textbf{Code block} (longer snippets with syntax highlighting):

\begin{verbatim}
\begin{lstlisting}[language=Python, caption={A simple neural network
training loop.}, label={lst:training}]
for epoch in range(num_epochs):
    model.train()
    optimizer.zero_grad()
    outputs = model(inputs)
    loss = criterion(outputs, labels)
    loss.backward()
    optimizer.step()
\end{lstlisting}
\end{verbatim}

\noindent
Supported languages include: \texttt{Python}, \texttt{Java},
\texttt{C}, \texttt{C++}, \texttt{JavaScript}, \texttt{SQL},
\texttt{bash}, \texttt{HTML}, \texttt{XML}.

\noindent
\textbf{Reference a listing:}

\begin{verbatim}
Listing~\ref{lst:training} shows...
\end{verbatim}


% ============================================================
\section*{7. Cross-References}
% ============================================================

\noindent
LaTeX can automatically number and reference almost anything —
chapters, sections, figures, tables, equations, and listings.
The pattern is always: label it with \verb|\label{}|, then
reference it with \verb|\ref{}|.

\begin{verbatim}
% Label a chapter, section, figure, table, or equation:
\chapter{Introduction}
\label{chapter1}          % labels this chapter

\section{Background}
\label{sec:background}    % labels this section

% Reference them anywhere in the document:
Chapter~\ref{chapter1} introduces...
Section~\ref{sec:background} covers...
\end{verbatim}

\noindent
\textbf{Naming convention for labels} (recommended):
\begin{itemize}
    \item Chapters: \texttt{chapter1}, \texttt{chapter2}
    \item Sections: \texttt{sec:background}, \texttt{sec:results}
    \item Figures:  \texttt{fig:architecture}, \texttt{fig:results}
    \item Tables:   \texttt{tab:comparison}, \texttt{tab:results}
    \item Equations:\texttt{eq:accuracy}, \texttt{eq:loss}
    \item Listings: \texttt{lst:training}, \texttt{lst:algorithm}
\end{itemize}


% ============================================================
\section*{8. Citations}
% ============================================================

\noindent
This template uses ACM numbered citations. Add your sources to
\texttt{References.bib}, then cite them in your text.

\begin{verbatim}
% Single citation — produces [1]
\cite{smith2024}

% Multiple citations — produces [1, 2, 3]
\cite{smith2024, jones2023, lee2022}

% Named author in prose — you write the name, LaTeX adds [1]
Smith et al. \cite{smith2024} found that...
Smith and Jones \cite{smithjones2024} proposed...
\end{verbatim}

\noindent
The numbers are assigned automatically in the order citations
first appear in your text. You do not need to number them manually.


% ============================================================
\section*{9. Footnotes}
% ============================================================

\begin{verbatim}
This is a sentence with a footnote.\footnote{The footnote text
appears at the bottom of the page automatically.}
\end{verbatim}

\noindent
Use footnotes sparingly in a thesis. Prefer in-text explanation
or a citation over a footnote where possible.


% ============================================================
\section*{10. Abbreviations (Glossary)}
% ============================================================

\noindent
Abbreviations are defined in \texttt{Lists.tex} and used in
your chapter text as follows:

\begin{verbatim}
% In Lists.tex:
\newacronym{ml}{ML}{Machine Learning}

% In your chapter:
\gls{ml} models have become widely used...
% First use renders as: "Machine Learning (ML) models..."
% Subsequent uses render as: "ML models..."
\end{verbatim}

\noindent
If you do not want the automatic expansion behavior, you can
also just write the abbreviation directly in your text without
using \verb|\gls{}|. The abbreviation list in the front matter
will still be generated from your \texttt{Lists.tex} entries.


% ============================================================
\section*{11. Common Mistakes}
% ============================================================

\begin{itemize}
    \item \textbf{Special characters}: The following characters
    have special meaning in LaTeX and must be escaped:
    \begin{center}
    \begin{tabular}{ll}
        \texttt{\%} → \texttt{\textbackslash\%} &
        \texttt{\$} → \texttt{\textbackslash\$} \\
        \texttt{\&} → \texttt{\textbackslash\&} &
        \texttt{\_} → \texttt{\textbackslash\_} \\
        \texttt{\#} → \texttt{\textbackslash\#} &
        \texttt{\{} → \texttt{\textbackslash\{} \\
    \end{tabular}
    \end{center}

    \item \textbf{Quotation marks}: Use \verb|``text''| not
    \verb|"text"| — straight quotes do not typeset correctly.

    \item \textbf{Figure placement}: Use \texttt{[H]} (from the
    \texttt{float} package) to place a figure exactly where you
    put it in the source. Without \texttt{[H]}, LaTeX may move
    figures to the top or bottom of a page.

    \item \textbf{References not updating}: Run pdflatex, then
    bibtex, then pdflatex twice. In Overleaf, just click
    Recompile — it handles this automatically.

    \item \textbf{Undefined references}: If you see \textbf{??}
    instead of a number, your \verb|\label{}| key does not match
    your \verb|\ref{}| key. Check spelling and capitalization.

    \item \textbf{Missing citation}: If \verb|\cite{KEY}| shows
    \textbf{[?]}, the key does not exist in \texttt{References.bib}
    or BibTeX has not been run yet. Check your bib file.
\end{itemize}


% ============================================================
\section*{12. Useful Resources}
% ============================================================

\begin{itemize}
    \item \textbf{Overleaf documentation}: \url{https://www.overleaf.com/learn}
    \item \textbf{LaTeX symbol lookup}: \url{https://detexify.kirelabs.org}
          (draw a symbol, get the LaTeX command)
    \item \textbf{Table generator}: \url{https://www.tablesgenerator.com}
          (build tables visually, export LaTeX code)
    \item \textbf{BibTeX format guide}:
          \url{https://www.acm.org/publications/authors/reference-formatting}
    \item \textbf{Math reference}: \url{https://en.wikibooks.org/wiki/LaTeX/Mathematics}
\end{itemize}

%   Remove it before your final submission.
% ============================================================

\chapter*{LaTeX Quick Reference}
\addcontentsline{toc}{chapter}{LaTeX Quick Reference}

\noindent
This chapter is a reference for common LaTeX features used in
CS theses. It is intended for students who are new to LaTeX or
who need a quick reminder of the syntax for specific tasks.

% ============================================================
\section*{1. Basic Text Formatting}
% ============================================================

\noindent
LaTeX handles most formatting automatically. A few things you
will use frequently:

\begin{itemize}
    \item \textbf{Bold text}: \verb|\textbf{your text}|
    \item \textit{Italic text}: \verb|\textit{your text}|
    \item \texttt{Code inline}: \verb|\texttt{your text}| or \verb|\verb|CODE||
    \item Line break within a paragraph: \verb|\\|
    \item New paragraph: leave a blank line in the source
    \item Em dash (—): \verb|---| \quad En dash (–): \verb|--|
    \item Quotation marks: \verb|``text''| renders as ``text''
\end{itemize}


% ============================================================
\section*{2. Lists}
% ============================================================

\noindent
\textbf{Bulleted list:}

\begin{verbatim}
\begin{itemize}
    \item First item
    \item Second item
    \item Third item
\end{itemize}
\end{verbatim}

Renders as:
\begin{itemize}
    \item First item
    \item Second item
    \item Third item
\end{itemize}

\noindent
\textbf{Numbered list:}

\begin{verbatim}
\begin{enumerate}
    \item First item
    \item Second item
\end{enumerate}
\end{verbatim}

\noindent
\textbf{Custom labels} (for research questions, contributions, etc.):

\begin{verbatim}
\begin{enumerate}
    \item[\textbf{RQ1}] First research question.
    \item[\textbf{RQ2}] Second research question.
\end{enumerate}
\end{verbatim}


% ============================================================
\section*{3. Figures}
% ============================================================

\noindent
\textbf{Single figure:}

\begin{verbatim}
\begin{figure}[H]
    \centering
    \includegraphics[width=0.7\textwidth]{Images/your_figure.png}
    \caption{A descriptive caption explaining the figure.}
    \label{fig:your_label}
\end{figure}
\end{verbatim}

\noindent
Key options for \verb|\includegraphics|:
\begin{itemize}
    \item \verb|width=0.7\textwidth| — 70\% of text width (adjust as needed)
    \item \verb|width=\textwidth| — full text width
    \item \verb|height=5cm| — fixed height (width scales automatically)
    \item \verb|scale=0.5| — scale to 50\% of original size
\end{itemize}

\noindent
\textbf{Reference a figure in your text:}

\begin{verbatim}
Figure~\ref{fig:your_label} shows...
As shown in Figure~\ref{fig:your_label},...
\end{verbatim}

\noindent
The \verb|~| prevents a line break between ``Figure'' and the number.

\noindent
\textbf{Side-by-side figures:}

\begin{verbatim}
\begin{figure}[H]
    \centering
    \begin{subfigure}[t]{0.48\textwidth}
        \centering
        \includegraphics[width=\textwidth]{Images/figure_a.png}
        \caption{Caption for the left figure.}
        \label{fig:left}
    \end{subfigure}
    \hfill
    \begin{subfigure}[t]{0.48\textwidth}
        \centering
        \includegraphics[width=\textwidth]{Images/figure_b.png}
        \caption{Caption for the right figure.}
        \label{fig:right}
    \end{subfigure}
    \caption{Overall caption for both figures.}
    \label{fig:both}
\end{figure}
\end{verbatim}

\noindent
\textbf{Important:} Put all image files in the \texttt{Images/} folder.
Supported formats: PNG, JPG, PDF. PNG is recommended for screenshots
and diagrams. PDF is recommended for vector graphics.


% ============================================================
\section*{4. Tables}
% ============================================================

\noindent
\textbf{Basic table:}

\begin{verbatim}
\begin{table}[H]
    \centering
    \caption{Caption goes ABOVE the table.}
    \label{tab:your_label}
    \begin{tabular}{lcc}
        \toprule
        \textbf{Column 1} & \textbf{Column 2} & \textbf{Column 3} \\
        \midrule
        Row 1 data        & 0.92              & 0.88              \\
        Row 2 data        & 0.87              & 0.91              \\
        Row 3 data        & 0.95              & 0.93              \\
        \bottomrule
    \end{tabular}
\end{table}
\end{verbatim}

\noindent
Column alignment characters:
\begin{itemize}
    \item \texttt{l} — left aligned
    \item \texttt{c} — centered
    \item \texttt{r} — right aligned
    \item \texttt{p\{3cm\}} — fixed-width column with text wrapping
\end{itemize}

\noindent
\textbf{Reference a table:}

\begin{verbatim}
Table~\ref{tab:your_label} shows...
\end{verbatim}

\noindent
\textbf{Auto-width table} (fills text width, good for many columns):

\begin{verbatim}
\begin{table}[H]
    \centering
    \caption{Caption here.}
    \label{tab:wide}
    \begin{tabularx}{\textwidth}{lXX}
        \toprule
        \textbf{Method} & \textbf{Description} & \textbf{Result} \\
        \midrule
        Method A & Describe it here & 92.3\% \\
        Method B & Describe it here & 88.7\% \\
        \bottomrule
    \end{tabularx}
\end{table}
\end{verbatim}

\noindent
\texttt{X} columns expand to fill available space equally.


% ============================================================
\section*{5. Equations}
% ============================================================

\noindent
\textbf{Numbered equation:}

\begin{verbatim}
\begin{equation}
    \text{Accuracy} = \frac{TP + TN}{TP + TN + FP + FN}
    \label{eq:accuracy}
\end{equation}
\end{verbatim}

\noindent
Renders as:
\begin{equation}
    \text{Accuracy} = \frac{TP + TN}{TP + TN + FP + FN}
    \label{eq:accuracy}
\end{equation}

\noindent
\textbf{Reference it in text:}

\begin{verbatim}
Equation~\ref{eq:accuracy} defines...
\end{verbatim}

\noindent
\textbf{Inline math} (within a sentence):

\begin{verbatim}
The learning rate $\alpha = 0.001$ was fixed throughout training.
\end{verbatim}

\noindent
\textbf{Unnumbered equation:}

\begin{verbatim}
\[
    y = mx + b
\]
\end{verbatim}

\noindent
\textbf{Multi-line aligned equations:}

\begin{verbatim}
\begin{align}
    f(x) &= x^2 + 2x + 1 \\
         &= (x + 1)^2
    \label{eq:multi}
\end{align}
\end{verbatim}

\noindent
Common math symbols:
\begin{itemize}
    \item Fraction: \verb|\frac{numerator}{denominator}|
    \item Subscript: \verb|x_i| → $x_i$
    \item Superscript: \verb|x^2| → $x^2$
    \item Greek: \verb|\alpha \beta \gamma \sigma \mu| → $\alpha\ \beta\ \gamma\ \sigma\ \mu$
    \item Sum: \verb|\sum_{i=1}^{n}| → $\sum_{i=1}^{n}$
    \item Square root: \verb|\sqrt{x}| → $\sqrt{x}$
    \item Infinity: \verb|\infty| → $\infty$
    \item Approximately: \verb|\approx| → $\approx$
    \item Less/greater or equal: \verb|\leq \geq| → $\leq\ \geq$
\end{itemize}


% ============================================================
\section*{6. Code Listings}
% ============================================================

\noindent
\textbf{Inline code} (short snippets in a sentence):

\begin{verbatim}
The function \texttt{train\_model()} accepts a learning rate parameter.
\end{verbatim}

\noindent
\textbf{Code block} (longer snippets with syntax highlighting):

\begin{verbatim}
\begin{lstlisting}[language=Python, caption={A simple neural network
training loop.}, label={lst:training}]
for epoch in range(num_epochs):
    model.train()
    optimizer.zero_grad()
    outputs = model(inputs)
    loss = criterion(outputs, labels)
    loss.backward()
    optimizer.step()
\end{lstlisting}
\end{verbatim}

\noindent
Supported languages include: \texttt{Python}, \texttt{Java},
\texttt{C}, \texttt{C++}, \texttt{JavaScript}, \texttt{SQL},
\texttt{bash}, \texttt{HTML}, \texttt{XML}.

\noindent
\textbf{Reference a listing:}

\begin{verbatim}
Listing~\ref{lst:training} shows...
\end{verbatim}


% ============================================================
\section*{7. Cross-References}
% ============================================================

\noindent
LaTeX can automatically number and reference almost anything —
chapters, sections, figures, tables, equations, and listings.
The pattern is always: label it with \verb|\label{}|, then
reference it with \verb|\ref{}|.

\begin{verbatim}
% Label a chapter, section, figure, table, or equation:
\chapter{Introduction}
\label{chapter1}          % labels this chapter

\section{Background}
\label{sec:background}    % labels this section

% Reference them anywhere in the document:
Chapter~\ref{chapter1} introduces...
Section~\ref{sec:background} covers...
\end{verbatim}

\noindent
\textbf{Naming convention for labels} (recommended):
\begin{itemize}
    \item Chapters: \texttt{chapter1}, \texttt{chapter2}
    \item Sections: \texttt{sec:background}, \texttt{sec:results}
    \item Figures:  \texttt{fig:architecture}, \texttt{fig:results}
    \item Tables:   \texttt{tab:comparison}, \texttt{tab:results}
    \item Equations:\texttt{eq:accuracy}, \texttt{eq:loss}
    \item Listings: \texttt{lst:training}, \texttt{lst:algorithm}
\end{itemize}


% ============================================================
\section*{8. Citations}
% ============================================================

\noindent
This template uses ACM numbered citations. Add your sources to
\texttt{References.bib}, then cite them in your text.

\begin{verbatim}
% Single citation — produces [1]
\cite{smith2024}

% Multiple citations — produces [1, 2, 3]
\cite{smith2024, jones2023, lee2022}

% Named author in prose — you write the name, LaTeX adds [1]
Smith et al. \cite{smith2024} found that...
Smith and Jones \cite{smithjones2024} proposed...
\end{verbatim}

\noindent
The numbers are assigned automatically in the order citations
first appear in your text. You do not need to number them manually.


% ============================================================
\section*{9. Footnotes}
% ============================================================

\begin{verbatim}
This is a sentence with a footnote.\footnote{The footnote text
appears at the bottom of the page automatically.}
\end{verbatim}

\noindent
Use footnotes sparingly in a thesis. Prefer in-text explanation
or a citation over a footnote where possible.


% ============================================================
\section*{10. Abbreviations (Glossary)}
% ============================================================

\noindent
Abbreviations are defined in \texttt{Lists.tex} and used in
your chapter text as follows:

\begin{verbatim}
% In Lists.tex:
\newacronym{ml}{ML}{Machine Learning}

% In your chapter:
\gls{ml} models have become widely used...
% First use renders as: "Machine Learning (ML) models..."
% Subsequent uses render as: "ML models..."
\end{verbatim}

\noindent
If you do not want the automatic expansion behavior, you can
also just write the abbreviation directly in your text without
using \verb|\gls{}|. The abbreviation list in the front matter
will still be generated from your \texttt{Lists.tex} entries.


% ============================================================
\section*{11. Common Mistakes}
% ============================================================

\begin{itemize}
    \item \textbf{Special characters}: The following characters
    have special meaning in LaTeX and must be escaped:
    \begin{center}
    \begin{tabular}{ll}
        \texttt{\%} → \texttt{\textbackslash\%} &
        \texttt{\$} → \texttt{\textbackslash\$} \\
        \texttt{\&} → \texttt{\textbackslash\&} &
        \texttt{\_} → \texttt{\textbackslash\_} \\
        \texttt{\#} → \texttt{\textbackslash\#} &
        \texttt{\{} → \texttt{\textbackslash\{} \\
    \end{tabular}
    \end{center}

    \item \textbf{Quotation marks}: Use \verb|``text''| not
    \verb|"text"| — straight quotes do not typeset correctly.

    \item \textbf{Figure placement}: Use \texttt{[H]} (from the
    \texttt{float} package) to place a figure exactly where you
    put it in the source. Without \texttt{[H]}, LaTeX may move
    figures to the top or bottom of a page.

    \item \textbf{References not updating}: Run pdflatex, then
    bibtex, then pdflatex twice. In Overleaf, just click
    Recompile — it handles this automatically.

    \item \textbf{Undefined references}: If you see \textbf{??}
    instead of a number, your \verb|\label{}| key does not match
    your \verb|\ref{}| key. Check spelling and capitalization.

    \item \textbf{Missing citation}: If \verb|\cite{KEY}| shows
    \textbf{[?]}, the key does not exist in \texttt{References.bib}
    or BibTeX has not been run yet. Check your bib file.
\end{itemize}


% ============================================================
\section*{12. Useful Resources}
% ============================================================

\begin{itemize}
    \item \textbf{Overleaf documentation}: \url{https://www.overleaf.com/learn}
    \item \textbf{LaTeX symbol lookup}: \url{https://detexify.kirelabs.org}
          (draw a symbol, get the LaTeX command)
    \item \textbf{Table generator}: \url{https://www.tablesgenerator.com}
          (build tables visually, export LaTeX code)
    \item \textbf{BibTeX format guide}:
          \url{https://www.acm.org/publications/authors/reference-formatting}
    \item \textbf{Math reference}: \url{https://en.wikibooks.org/wiki/LaTeX/Mathematics}
\end{itemize}

%   Remove it before your final submission.
% ============================================================

\chapter*{LaTeX Quick Reference}
\addcontentsline{toc}{chapter}{LaTeX Quick Reference}

\noindent
This chapter is a reference for common LaTeX features used in
CS theses. It is intended for students who are new to LaTeX or
who need a quick reminder of the syntax for specific tasks.

% ============================================================
\section*{1. Basic Text Formatting}
% ============================================================

\noindent
LaTeX handles most formatting automatically. A few things you
will use frequently:

\begin{itemize}
    \item \textbf{Bold text}: \verb|\textbf{your text}|
    \item \textit{Italic text}: \verb|\textit{your text}|
    \item \texttt{Code inline}: \verb|\texttt{your text}| or \verb|\verb|CODE||
    \item Line break within a paragraph: \verb|\\|
    \item New paragraph: leave a blank line in the source
    \item Em dash (—): \verb|---| \quad En dash (–): \verb|--|
    \item Quotation marks: \verb|``text''| renders as ``text''
\end{itemize}


% ============================================================
\section*{2. Lists}
% ============================================================

\noindent
\textbf{Bulleted list:}

\begin{verbatim}
\begin{itemize}
    \item First item
    \item Second item
    \item Third item
\end{itemize}
\end{verbatim}

Renders as:
\begin{itemize}
    \item First item
    \item Second item
    \item Third item
\end{itemize}

\noindent
\textbf{Numbered list:}

\begin{verbatim}
\begin{enumerate}
    \item First item
    \item Second item
\end{enumerate}
\end{verbatim}

\noindent
\textbf{Custom labels} (for research questions, contributions, etc.):

\begin{verbatim}
\begin{enumerate}
    \item[\textbf{RQ1}] First research question.
    \item[\textbf{RQ2}] Second research question.
\end{enumerate}
\end{verbatim}


% ============================================================
\section*{3. Figures}
% ============================================================

\noindent
\textbf{Single figure:}

\begin{verbatim}
\begin{figure}[H]
    \centering
    \includegraphics[width=0.7\textwidth]{Images/your_figure.png}
    \caption{A descriptive caption explaining the figure.}
    \label{fig:your_label}
\end{figure}
\end{verbatim}

\noindent
Key options for \verb|\includegraphics|:
\begin{itemize}
    \item \verb|width=0.7\textwidth| — 70\% of text width (adjust as needed)
    \item \verb|width=\textwidth| — full text width
    \item \verb|height=5cm| — fixed height (width scales automatically)
    \item \verb|scale=0.5| — scale to 50\% of original size
\end{itemize}

\noindent
\textbf{Reference a figure in your text:}

\begin{verbatim}
Figure~\ref{fig:your_label} shows...
As shown in Figure~\ref{fig:your_label},...
\end{verbatim}

\noindent
The \verb|~| prevents a line break between ``Figure'' and the number.

\noindent
\textbf{Side-by-side figures:}

\begin{verbatim}
\begin{figure}[H]
    \centering
    \begin{subfigure}[t]{0.48\textwidth}
        \centering
        \includegraphics[width=\textwidth]{Images/figure_a.png}
        \caption{Caption for the left figure.}
        \label{fig:left}
    \end{subfigure}
    \hfill
    \begin{subfigure}[t]{0.48\textwidth}
        \centering
        \includegraphics[width=\textwidth]{Images/figure_b.png}
        \caption{Caption for the right figure.}
        \label{fig:right}
    \end{subfigure}
    \caption{Overall caption for both figures.}
    \label{fig:both}
\end{figure}
\end{verbatim}

\noindent
\textbf{Important:} Put all image files in the \texttt{Images/} folder.
Supported formats: PNG, JPG, PDF. PNG is recommended for screenshots
and diagrams. PDF is recommended for vector graphics.


% ============================================================
\section*{4. Tables}
% ============================================================

\noindent
\textbf{Basic table:}

\begin{verbatim}
\begin{table}[H]
    \centering
    \caption{Caption goes ABOVE the table.}
    \label{tab:your_label}
    \begin{tabular}{lcc}
        \toprule
        \textbf{Column 1} & \textbf{Column 2} & \textbf{Column 3} \\
        \midrule
        Row 1 data        & 0.92              & 0.88              \\
        Row 2 data        & 0.87              & 0.91              \\
        Row 3 data        & 0.95              & 0.93              \\
        \bottomrule
    \end{tabular}
\end{table}
\end{verbatim}

\noindent
Column alignment characters:
\begin{itemize}
    \item \texttt{l} — left aligned
    \item \texttt{c} — centered
    \item \texttt{r} — right aligned
    \item \texttt{p\{3cm\}} — fixed-width column with text wrapping
\end{itemize}

\noindent
\textbf{Reference a table:}

\begin{verbatim}
Table~\ref{tab:your_label} shows...
\end{verbatim}

\noindent
\textbf{Auto-width table} (fills text width, good for many columns):

\begin{verbatim}
\begin{table}[H]
    \centering
    \caption{Caption here.}
    \label{tab:wide}
    \begin{tabularx}{\textwidth}{lXX}
        \toprule
        \textbf{Method} & \textbf{Description} & \textbf{Result} \\
        \midrule
        Method A & Describe it here & 92.3\% \\
        Method B & Describe it here & 88.7\% \\
        \bottomrule
    \end{tabularx}
\end{table}
\end{verbatim}

\noindent
\texttt{X} columns expand to fill available space equally.


% ============================================================
\section*{5. Equations}
% ============================================================

\noindent
\textbf{Numbered equation:}

\begin{verbatim}
\begin{equation}
    \text{Accuracy} = \frac{TP + TN}{TP + TN + FP + FN}
    \label{eq:accuracy}
\end{equation}
\end{verbatim}

\noindent
Renders as:
\begin{equation}
    \text{Accuracy} = \frac{TP + TN}{TP + TN + FP + FN}
    \label{eq:accuracy}
\end{equation}

\noindent
\textbf{Reference it in text:}

\begin{verbatim}
Equation~\ref{eq:accuracy} defines...
\end{verbatim}

\noindent
\textbf{Inline math} (within a sentence):

\begin{verbatim}
The learning rate $\alpha = 0.001$ was fixed throughout training.
\end{verbatim}

\noindent
\textbf{Unnumbered equation:}

\begin{verbatim}
\[
    y = mx + b
\]
\end{verbatim}

\noindent
\textbf{Multi-line aligned equations:}

\begin{verbatim}
\begin{align}
    f(x) &= x^2 + 2x + 1 \\
         &= (x + 1)^2
    \label{eq:multi}
\end{align}
\end{verbatim}

\noindent
Common math symbols:
\begin{itemize}
    \item Fraction: \verb|\frac{numerator}{denominator}|
    \item Subscript: \verb|x_i| → $x_i$
    \item Superscript: \verb|x^2| → $x^2$
    \item Greek: \verb|\alpha \beta \gamma \sigma \mu| → $\alpha\ \beta\ \gamma\ \sigma\ \mu$
    \item Sum: \verb|\sum_{i=1}^{n}| → $\sum_{i=1}^{n}$
    \item Square root: \verb|\sqrt{x}| → $\sqrt{x}$
    \item Infinity: \verb|\infty| → $\infty$
    \item Approximately: \verb|\approx| → $\approx$
    \item Less/greater or equal: \verb|\leq \geq| → $\leq\ \geq$
\end{itemize}


% ============================================================
\section*{6. Code Listings}
% ============================================================

\noindent
\textbf{Inline code} (short snippets in a sentence):

\begin{verbatim}
The function \texttt{train\_model()} accepts a learning rate parameter.
\end{verbatim}

\noindent
\textbf{Code block} (longer snippets with syntax highlighting):

\begin{verbatim}
\begin{lstlisting}[language=Python, caption={A simple neural network
training loop.}, label={lst:training}]
for epoch in range(num_epochs):
    model.train()
    optimizer.zero_grad()
    outputs = model(inputs)
    loss = criterion(outputs, labels)
    loss.backward()
    optimizer.step()
\end{lstlisting}
\end{verbatim}

\noindent
Supported languages include: \texttt{Python}, \texttt{Java},
\texttt{C}, \texttt{C++}, \texttt{JavaScript}, \texttt{SQL},
\texttt{bash}, \texttt{HTML}, \texttt{XML}.

\noindent
\textbf{Reference a listing:}

\begin{verbatim}
Listing~\ref{lst:training} shows...
\end{verbatim}


% ============================================================
\section*{7. Cross-References}
% ============================================================

\noindent
LaTeX can automatically number and reference almost anything —
chapters, sections, figures, tables, equations, and listings.
The pattern is always: label it with \verb|\label{}|, then
reference it with \verb|\ref{}|.

\begin{verbatim}
% Label a chapter, section, figure, table, or equation:
\chapter{Introduction}
\label{chapter1}          % labels this chapter

\section{Background}
\label{sec:background}    % labels this section

% Reference them anywhere in the document:
Chapter~\ref{chapter1} introduces...
Section~\ref{sec:background} covers...
\end{verbatim}

\noindent
\textbf{Naming convention for labels} (recommended):
\begin{itemize}
    \item Chapters: \texttt{chapter1}, \texttt{chapter2}
    \item Sections: \texttt{sec:background}, \texttt{sec:results}
    \item Figures:  \texttt{fig:architecture}, \texttt{fig:results}
    \item Tables:   \texttt{tab:comparison}, \texttt{tab:results}
    \item Equations:\texttt{eq:accuracy}, \texttt{eq:loss}
    \item Listings: \texttt{lst:training}, \texttt{lst:algorithm}
\end{itemize}


% ============================================================
\section*{8. Citations}
% ============================================================

\noindent
This template uses ACM numbered citations. Add your sources to
\texttt{References.bib}, then cite them in your text.

\begin{verbatim}
% Single citation — produces [1]
\cite{smith2024}

% Multiple citations — produces [1, 2, 3]
\cite{smith2024, jones2023, lee2022}

% Named author in prose — you write the name, LaTeX adds [1]
Smith et al. \cite{smith2024} found that...
Smith and Jones \cite{smithjones2024} proposed...
\end{verbatim}

\noindent
The numbers are assigned automatically in the order citations
first appear in your text. You do not need to number them manually.


% ============================================================
\section*{9. Footnotes}
% ============================================================

\begin{verbatim}
This is a sentence with a footnote.\footnote{The footnote text
appears at the bottom of the page automatically.}
\end{verbatim}

\noindent
Use footnotes sparingly in a thesis. Prefer in-text explanation
or a citation over a footnote where possible.


% ============================================================
\section*{10. Abbreviations (Glossary)}
% ============================================================

\noindent
Abbreviations are defined in \texttt{Lists.tex} and used in
your chapter text as follows:

\begin{verbatim}
% In Lists.tex:
\newacronym{ml}{ML}{Machine Learning}

% In your chapter:
\gls{ml} models have become widely used...
% First use renders as: "Machine Learning (ML) models..."
% Subsequent uses render as: "ML models..."
\end{verbatim}

\noindent
If you do not want the automatic expansion behavior, you can
also just write the abbreviation directly in your text without
using \verb|\gls{}|. The abbreviation list in the front matter
will still be generated from your \texttt{Lists.tex} entries.


% ============================================================
\section*{11. Common Mistakes}
% ============================================================

\begin{itemize}
    \item \textbf{Special characters}: The following characters
    have special meaning in LaTeX and must be escaped:
    \begin{center}
    \begin{tabular}{ll}
        \texttt{\%} → \texttt{\textbackslash\%} &
        \texttt{\$} → \texttt{\textbackslash\$} \\
        \texttt{\&} → \texttt{\textbackslash\&} &
        \texttt{\_} → \texttt{\textbackslash\_} \\
        \texttt{\#} → \texttt{\textbackslash\#} &
        \texttt{\{} → \texttt{\textbackslash\{} \\
    \end{tabular}
    \end{center}

    \item \textbf{Quotation marks}: Use \verb|``text''| not
    \verb|"text"| — straight quotes do not typeset correctly.

    \item \textbf{Figure placement}: Use \texttt{[H]} (from the
    \texttt{float} package) to place a figure exactly where you
    put it in the source. Without \texttt{[H]}, LaTeX may move
    figures to the top or bottom of a page.

    \item \textbf{References not updating}: Run pdflatex, then
    bibtex, then pdflatex twice. In Overleaf, just click
    Recompile — it handles this automatically.

    \item \textbf{Undefined references}: If you see \textbf{??}
    instead of a number, your \verb|\label{}| key does not match
    your \verb|\ref{}| key. Check spelling and capitalization.

    \item \textbf{Missing citation}: If \verb|\cite{KEY}| shows
    \textbf{[?]}, the key does not exist in \texttt{References.bib}
    or BibTeX has not been run yet. Check your bib file.
\end{itemize}


% ============================================================
\section*{12. Useful Resources}
% ============================================================

\begin{itemize}
    \item \textbf{Overleaf documentation}: \url{https://www.overleaf.com/learn}
    \item \textbf{LaTeX symbol lookup}: \url{https://detexify.kirelabs.org}
          (draw a symbol, get the LaTeX command)
    \item \textbf{Table generator}: \url{https://www.tablesgenerator.com}
          (build tables visually, export LaTeX code)
    \item \textbf{BibTeX format guide}:
          \url{https://www.acm.org/publications/authors/reference-formatting}
    \item \textbf{Math reference}: \url{https://en.wikibooks.org/wiki/LaTeX/Mathematics}
\end{itemize}

%   Remove it before your final submission.
% ============================================================

\chapter*{LaTeX Quick Reference}
\addcontentsline{toc}{chapter}{LaTeX Quick Reference}

\noindent
This chapter is a reference for common LaTeX features used in
CS theses. It is intended for students who are new to LaTeX or
who need a quick reminder of the syntax for specific tasks.

% ============================================================
\section*{1. Basic Text Formatting}
% ============================================================

\noindent
LaTeX handles most formatting automatically. A few things you
will use frequently:

\begin{itemize}
    \item \textbf{Bold text}: \verb|\textbf{your text}|
    \item \textit{Italic text}: \verb|\textit{your text}|
    \item \texttt{Code inline}: \verb|\texttt{your text}| or \verb|\verb|CODE||
    \item Line break within a paragraph: \verb|\\|
    \item New paragraph: leave a blank line in the source
    \item Em dash (—): \verb|---| \quad En dash (–): \verb|--|
    \item Quotation marks: \verb|``text''| renders as ``text''
\end{itemize}


% ============================================================
\section*{2. Lists}
% ============================================================

\noindent
\textbf{Bulleted list:}

\begin{verbatim}
\begin{itemize}
    \item First item
    \item Second item
    \item Third item
\end{itemize}
\end{verbatim}

Renders as:
\begin{itemize}
    \item First item
    \item Second item
    \item Third item
\end{itemize}

\noindent
\textbf{Numbered list:}

\begin{verbatim}
\begin{enumerate}
    \item First item
    \item Second item
\end{enumerate}
\end{verbatim}

\noindent
\textbf{Custom labels} (for research questions, contributions, etc.):

\begin{verbatim}
\begin{enumerate}
    \item[\textbf{RQ1}] First research question.
    \item[\textbf{RQ2}] Second research question.
\end{enumerate}
\end{verbatim}


% ============================================================
\section*{3. Figures}
% ============================================================

\noindent
\textbf{Single figure:}

\begin{verbatim}
\begin{figure}[H]
    \centering
    \includegraphics[width=0.7\textwidth]{Images/your_figure.png}
    \caption{A descriptive caption explaining the figure.}
    \label{fig:your_label}
\end{figure}
\end{verbatim}

\noindent
Key options for \verb|\includegraphics|:
\begin{itemize}
    \item \verb|width=0.7\textwidth| — 70\% of text width (adjust as needed)
    \item \verb|width=\textwidth| — full text width
    \item \verb|height=5cm| — fixed height (width scales automatically)
    \item \verb|scale=0.5| — scale to 50\% of original size
\end{itemize}

\noindent
\textbf{Reference a figure in your text:}

\begin{verbatim}
Figure~\ref{fig:your_label} shows...
As shown in Figure~\ref{fig:your_label},...
\end{verbatim}

\noindent
The \verb|~| prevents a line break between ``Figure'' and the number.

\noindent
\textbf{Side-by-side figures:}

\begin{verbatim}
\begin{figure}[H]
    \centering
    \begin{subfigure}[t]{0.48\textwidth}
        \centering
        \includegraphics[width=\textwidth]{Images/figure_a.png}
        \caption{Caption for the left figure.}
        \label{fig:left}
    \end{subfigure}
    \hfill
    \begin{subfigure}[t]{0.48\textwidth}
        \centering
        \includegraphics[width=\textwidth]{Images/figure_b.png}
        \caption{Caption for the right figure.}
        \label{fig:right}
    \end{subfigure}
    \caption{Overall caption for both figures.}
    \label{fig:both}
\end{figure}
\end{verbatim}

\noindent
\textbf{Important:} Put all image files in the \texttt{Images/} folder.
Supported formats: PNG, JPG, PDF. PNG is recommended for screenshots
and diagrams. PDF is recommended for vector graphics.


% ============================================================
\section*{4. Tables}
% ============================================================

\noindent
\textbf{Basic table:}

\begin{verbatim}
\begin{table}[H]
    \centering
    \caption{Caption goes ABOVE the table.}
    \label{tab:your_label}
    \begin{tabular}{lcc}
        \toprule
        \textbf{Column 1} & \textbf{Column 2} & \textbf{Column 3} \\
        \midrule
        Row 1 data        & 0.92              & 0.88              \\
        Row 2 data        & 0.87              & 0.91              \\
        Row 3 data        & 0.95              & 0.93              \\
        \bottomrule
    \end{tabular}
\end{table}
\end{verbatim}

\noindent
Column alignment characters:
\begin{itemize}
    \item \texttt{l} — left aligned
    \item \texttt{c} — centered
    \item \texttt{r} — right aligned
    \item \texttt{p\{3cm\}} — fixed-width column with text wrapping
\end{itemize}

\noindent
\textbf{Reference a table:}

\begin{verbatim}
Table~\ref{tab:your_label} shows...
\end{verbatim}

\noindent
\textbf{Auto-width table} (fills text width, good for many columns):

\begin{verbatim}
\begin{table}[H]
    \centering
    \caption{Caption here.}
    \label{tab:wide}
    \begin{tabularx}{\textwidth}{lXX}
        \toprule
        \textbf{Method} & \textbf{Description} & \textbf{Result} \\
        \midrule
        Method A & Describe it here & 92.3\% \\
        Method B & Describe it here & 88.7\% \\
        \bottomrule
    \end{tabularx}
\end{table}
\end{verbatim}

\noindent
\texttt{X} columns expand to fill available space equally.


% ============================================================
\section*{5. Equations}
% ============================================================

\noindent
\textbf{Numbered equation:}

\begin{verbatim}
\begin{equation}
    \text{Accuracy} = \frac{TP + TN}{TP + TN + FP + FN}
    \label{eq:accuracy}
\end{equation}
\end{verbatim}

\noindent
Renders as:
\begin{equation}
    \text{Accuracy} = \frac{TP + TN}{TP + TN + FP + FN}
    \label{eq:accuracy}
\end{equation}

\noindent
\textbf{Reference it in text:}

\begin{verbatim}
Equation~\ref{eq:accuracy} defines...
\end{verbatim}

\noindent
\textbf{Inline math} (within a sentence):

\begin{verbatim}
The learning rate $\alpha = 0.001$ was fixed throughout training.
\end{verbatim}

\noindent
\textbf{Unnumbered equation:}

\begin{verbatim}
\[
    y = mx + b
\]
\end{verbatim}

\noindent
\textbf{Multi-line aligned equations:}

\begin{verbatim}
\begin{align}
    f(x) &= x^2 + 2x + 1 \\
         &= (x + 1)^2
    \label{eq:multi}
\end{align}
\end{verbatim}

\noindent
Common math symbols:
\begin{itemize}
    \item Fraction: \verb|\frac{numerator}{denominator}|
    \item Subscript: \verb|x_i| → $x_i$
    \item Superscript: \verb|x^2| → $x^2$
    \item Greek: \verb|\alpha \beta \gamma \sigma \mu| → $\alpha\ \beta\ \gamma\ \sigma\ \mu$
    \item Sum: \verb|\sum_{i=1}^{n}| → $\sum_{i=1}^{n}$
    \item Square root: \verb|\sqrt{x}| → $\sqrt{x}$
    \item Infinity: \verb|\infty| → $\infty$
    \item Approximately: \verb|\approx| → $\approx$
    \item Less/greater or equal: \verb|\leq \geq| → $\leq\ \geq$
\end{itemize}


% ============================================================
\section*{6. Code Listings}
% ============================================================

\noindent
\textbf{Inline code} (short snippets in a sentence):

\begin{verbatim}
The function \texttt{train\_model()} accepts a learning rate parameter.
\end{verbatim}

\noindent
\textbf{Code block} (longer snippets with syntax highlighting):

\begin{verbatim}
\begin{lstlisting}[language=Python, caption={A simple neural network
training loop.}, label={lst:training}]
for epoch in range(num_epochs):
    model.train()
    optimizer.zero_grad()
    outputs = model(inputs)
    loss = criterion(outputs, labels)
    loss.backward()
    optimizer.step()
\end{lstlisting}
\end{verbatim}

\noindent
Supported languages include: \texttt{Python}, \texttt{Java},
\texttt{C}, \texttt{C++}, \texttt{JavaScript}, \texttt{SQL},
\texttt{bash}, \texttt{HTML}, \texttt{XML}.

\noindent
\textbf{Reference a listing:}

\begin{verbatim}
Listing~\ref{lst:training} shows...
\end{verbatim}


% ============================================================
\section*{7. Cross-References}
% ============================================================

\noindent
LaTeX can automatically number and reference almost anything —
chapters, sections, figures, tables, equations, and listings.
The pattern is always: label it with \verb|\label{}|, then
reference it with \verb|\ref{}|.

\begin{verbatim}
% Label a chapter, section, figure, table, or equation:
\chapter{Introduction}
\label{chapter1}          % labels this chapter

\section{Background}
\label{sec:background}    % labels this section

% Reference them anywhere in the document:
Chapter~\ref{chapter1} introduces...
Section~\ref{sec:background} covers...
\end{verbatim}

\noindent
\textbf{Naming convention for labels} (recommended):
\begin{itemize}
    \item Chapters: \texttt{chapter1}, \texttt{chapter2}
    \item Sections: \texttt{sec:background}, \texttt{sec:results}
    \item Figures:  \texttt{fig:architecture}, \texttt{fig:results}
    \item Tables:   \texttt{tab:comparison}, \texttt{tab:results}
    \item Equations:\texttt{eq:accuracy}, \texttt{eq:loss}
    \item Listings: \texttt{lst:training}, \texttt{lst:algorithm}
\end{itemize}


% ============================================================
\section*{8. Citations}
% ============================================================

\noindent
This template uses ACM numbered citations. Add your sources to
\texttt{References.bib}, then cite them in your text.

\begin{verbatim}
% Single citation — produces [1]
\cite{smith2024}

% Multiple citations — produces [1, 2, 3]
\cite{smith2024, jones2023, lee2022}

% Named author in prose — you write the name, LaTeX adds [1]
Smith et al. \cite{smith2024} found that...
Smith and Jones \cite{smithjones2024} proposed...
\end{verbatim}

\noindent
The numbers are assigned automatically in the order citations
first appear in your text. You do not need to number them manually.


% ============================================================
\section*{9. Footnotes}
% ============================================================

\begin{verbatim}
This is a sentence with a footnote.\footnote{The footnote text
appears at the bottom of the page automatically.}
\end{verbatim}

\noindent
Use footnotes sparingly in a thesis. Prefer in-text explanation
or a citation over a footnote where possible.


% ============================================================
\section*{10. Abbreviations (Glossary)}
% ============================================================

\noindent
Abbreviations are defined in \texttt{Lists.tex} and used in
your chapter text as follows:

\begin{verbatim}
% In Lists.tex:
\newacronym{ml}{ML}{Machine Learning}

% In your chapter:
\gls{ml} models have become widely used...
% First use renders as: "Machine Learning (ML) models..."
% Subsequent uses render as: "ML models..."
\end{verbatim}

\noindent
If you do not want the automatic expansion behavior, you can
also just write the abbreviation directly in your text without
using \verb|\gls{}|. The abbreviation list in the front matter
will still be generated from your \texttt{Lists.tex} entries.


% ============================================================
\section*{11. Common Mistakes}
% ============================================================

\begin{itemize}
    \item \textbf{Special characters}: The following characters
    have special meaning in LaTeX and must be escaped:
    \begin{center}
    \begin{tabular}{ll}
        \texttt{\%} → \texttt{\textbackslash\%} &
        \texttt{\$} → \texttt{\textbackslash\$} \\
        \texttt{\&} → \texttt{\textbackslash\&} &
        \texttt{\_} → \texttt{\textbackslash\_} \\
        \texttt{\#} → \texttt{\textbackslash\#} &
        \texttt{\{} → \texttt{\textbackslash\{} \\
    \end{tabular}
    \end{center}

    \item \textbf{Quotation marks}: Use \verb|``text''| not
    \verb|"text"| — straight quotes do not typeset correctly.

    \item \textbf{Figure placement}: Use \texttt{[H]} (from the
    \texttt{float} package) to place a figure exactly where you
    put it in the source. Without \texttt{[H]}, LaTeX may move
    figures to the top or bottom of a page.

    \item \textbf{References not updating}: Run pdflatex, then
    bibtex, then pdflatex twice. In Overleaf, just click
    Recompile — it handles this automatically.

    \item \textbf{Undefined references}: If you see \textbf{??}
    instead of a number, your \verb|\label{}| key does not match
    your \verb|\ref{}| key. Check spelling and capitalization.

    \item \textbf{Missing citation}: If \verb|\cite{KEY}| shows
    \textbf{[?]}, the key does not exist in \texttt{References.bib}
    or BibTeX has not been run yet. Check your bib file.
\end{itemize}


% ============================================================
\section*{12. Useful Resources}
% ============================================================

\begin{itemize}
    \item \textbf{Overleaf documentation}: \url{https://www.overleaf.com/learn}
    \item \textbf{LaTeX symbol lookup}: \url{https://detexify.kirelabs.org}
          (draw a symbol, get the LaTeX command)
    \item \textbf{Table generator}: \url{https://www.tablesgenerator.com}
          (build tables visually, export LaTeX code)
    \item \textbf{BibTeX format guide}:
          \url{https://www.acm.org/publications/authors/reference-formatting}
    \item \textbf{Math reference}: \url{https://en.wikibooks.org/wiki/LaTeX/Mathematics}
\end{itemize}
